\documentclass[10pt,a4paper]{amsart}
\usepackage[margin=19mm]{geometry}
\usepackage{amsmath,amssymb,mathtools}
\usepackage[scr=rsfs,cal=euler]{mathalfa}
\usepackage{xfrac}
\usepackage[unicode,hidelinks,bookmarks=false]{hyperref}
\newcommand{\ie}{i.e.\ }
\newcommand{\cf}{cf.\ }
\newcommand{\resp}{respectively\ }
\newcommand{\Subsection}{\subsection}
\providecommand{\nobreakdash}{\nobreak\hbox{-}}
\setlength{\emergencystretch}{2em}
\multlinegap0pt
\usepackage{bm}
\usepackage{bbm}
\usepackage{dsfont}
\usepackage{xspace}
\usepackage{cancel}
\usepackage{mathtools}
\usepackage{amsthm}
\newtheorem{theorem}{Theorem}
\newtheorem{lemma}[theorem]{Lemma}
\newcommand{\Q}{\mathbb{Q}}
\newcommand{\Z}{\mathbb{Z}}
\title[Euler characteristics of $SL(2,\mathbb{Z})$-orbit graphs of ]{Euler characteristics of $SL(2,\mathbb{Z})$-orbit graphs of origamis}
\author{Luke Jeffreys, Carlos Matheus}
\date{Source: arXiv:2602.21984v1 (CC BY 4.0)}
\begin{document}
\maketitle

\noindent\emph{Source and licence.} Euler characteristics of $SL(2,\mathbb{Z})$-orbit graphs of origamis. Version arXiv:2602.21984v1 (CC BY 4.0), distributed under the Creative Commons Attribution 4.0 International licence, as recorded in the arXiv licence metadata for this exact version and shown on the abstract page https://arxiv.org/abs/2602.21984v1. Full rights notice: licenses/arxiv-2602.21984-CC-BY-4.0.txt.\par\smallskip

\noindent\textit{Editorial scope.} Counted unit: the Lemma whose source label is \texttt{lem:TTZ3} in the article (source0.tex lines 892--899, source label \texttt{lem:TTZ3}), delivered together with the article's own complete original proof of it (source0.tex lines 901--931, one \texttt{proof} environment of 31 lines). No mathematics is written, repaired or re-derived: statement and proof are the article's own text line for line. Required context retained uncounted: none. Every definition, display and label of those lines is retained unchanged. Omitted from this package and not redistributed: 1--891, 900--900, 932--2020. Nothing inside a retained excerpt is omitted or altered: every retained body is delivered exactly as the article's lines are written, so no omission receipt of any kind accompanies this package. Citations: the retained excerpts carry no citation command at all, so no bibliography entry of the article is transcribed and no reference is redistributed. Reference adaptation: none was needed and none was made; every reference inside the delivered unit resolves inside it, so no unresolved reference or citation is produced. Printed slips kept literal: no printed word, symbol, hypothesis or number of the counted statement or of its proof was altered. Layout and class: the article's own class is \texttt{amsart} with packages including geometry, amsmath, amssymb, amscd, color, xcolor, mathrsfs, tikz, tikz-cd, float, comment, hyperref, hypcap, mathtools; the wrapper is the lane's pinned \texttt{amsart} editorial wrapper at 10pt on A4, which restates the theorem-like environments the retained text needs and adds the article's own macro definitions that the retained text needs (\texttt{\textbackslash Q}, \texttt{\textbackslash Z}), with extra wrapper packages (bm, bbm, dsfont, xspace, cancel, mathtools). The delivered numbering is the unit's own. Assets: no retained span contains a figure environment, a graphics-inclusion command, a PDF-page inclusion, a table or a TikZ picture; no image file is copied into this package, the package declares no asset dependency and no image file is redistributed with it. Licence: version arXiv:2602.21984v1 of this article is distributed under the Creative Commons Attribution 4.0 International licence, as recorded in the arXiv licence metadata for this exact version and shown on the abstract page https://arxiv.org/abs/2602.21984v1. A full rights notice is delivered at licenses/arxiv-2602.21984-CC-BY-4.0.txt. Characters: every retained excerpt is pure ASCII, so the delivered engine, XeLaTeX with its default Latin Modern text font, reports no missing character.
\begin{lemma}\label{lem:TTZ3}
    Let $D$ be a square discriminant. Then the matrices $A_{\sqrt{D}}(a,b,c)$ and $A_{\sqrt{D'}}(a',b',c')$ have the same eigenvectors if and only if
    \begin{itemize}
        \item[(i)] $D = m^2E$ and $D' = (m')^2E$ for some square discriminant $E$ and with $\gcd(m,m') = 1$; and 
        \item[(ii)] both $(a,b,c)$ and $(a',b',c')$ are integer multiples of a triple $(a_0,b_0,c_0)\in\Z^{3}$ with $a_0^{2}+b_0^{2}+c_0^{2} = E$.
    \end{itemize}
    In particular, the matrices $A_{\sqrt{D}}(a,b,c)$ and $A_{\sqrt{D}}(a',b',c')$ have the same eigenvectors if and only if $(a',b',c') = \pm(a,b,c)$ and we have $\omega(a,b,c)^{+} = \omega(-a,-b,-c)^{-}$ and $\omega(a,b,c)^{-} = \omega(-a,-b,-c)^{+}$.
\end{lemma}

\begin{proof}
    Suppose that $D = n^2$ and $D' = (n')^2$.
    
    If the matrices $A_{\sqrt{D}}(a,b,c)$ and $A_{\sqrt{D'}}(a',b',c')$ have the same eigenvectors then
    \begin{equation}\label{eqn:eigvals}
        \frac{a-ci}{b+n} = \frac{a'-c'i}{b'\pm n'}\,\,\,\,\text{and}\,\,\,\,\frac{a-ci}{b-n} = \frac{a'-c'i}{b'\mp n'}.
    \end{equation}
    This immediately forces $D = m^2E$ and $D' = (m')^2E$ for some square discriminant $E$ and we can force $\gcd(m,m') = 1$.

    Equating the real and imaginary parts of the first equality in Equation~\ref{eqn:eigvals} implies that
    \[a(b'\pm n') = a'(b+n)\,\,\,\,\,\text{and}\,\,\,\,\,c(b'\pm n')=c'(b+n).\]
    Multiplying by $c$ and $a$, respectively, and subtracting gives
    \[(a'c-ac')(b+n) = 0.\]

    If $b + n = 0$, then $a^2+b^2+c^2 = n^2$ forces $(a,b,c) = (0,-n,0)$ and also $(a',b',c') = (0,\pm n',0)$. Hence, both triples are integer multiples of $(a_0,b_0,c_0) = (0,\frac{n}{m},0)$ with $a_0^2+b_0^2+c_0^2 = E$.

    If $(a'c-ac') = \begin{vmatrix}
        a' & c' \\
        a & c
    \end{vmatrix} = 0$, then the vectors $(a,c)$ and $(a',c')$ are collinear. So $(a',c') = t(a,c)$ for some $t\in\Q$.
    
    If $a = c = 0$, then $(a,b,c) = (0,\pm n,0)$ and $(a',b',c') = (0,\pm n',0)$ and we are essentially in the previous situation again.
    
    So we may assume that one of $a$ or $c$ is non-zero. Hence, letting $t = \frac{a'}{a}$ or $\frac{c'}{c}$ appropriately, we have
    \[b'\pm n' = t(b+n) = tb+tn.\]
    Applying the same reasoning as above to the second equality in Equation~\ref{eqn:eigvals}, we will arrive at
    \[b'\mp n'= tb-tn.\]
    Therefore, we have $b' = tb$ and $n' = \pm tn$, giving $(a',b',c') = t(a,b,c) = tm(a_0,b_0,c_0)$ with $a_0^2+b_0^2+c_0^2 = E$.

    The converse follows immediately.
\end{proof}

\end{document}
